\section{Quantization of electromagnetic fields (10/30)}

\referto{\sakurai~7.6, \tong~5.3, 13.4.1, \scully~1.1}

\subsection{Classical Electromagnetic Field}

Let
\begin{equation}
 x^a=(ct,\bm r),
 \qquad
 A^a=(\phi/c,\bm A),
\end{equation}
and
\begin{equation}
 F_{ab}=\frac{\partial A_b}{\partial x^a}
 -\frac{\partial A_a}{\partial x^b}.
\end{equation}
The electromagnetic Lagrangian density is
\begin{equation}
 \mathcal L=-\frac{1}{4\mu}F^{ab}F_{ab}
 =\frac12\left(\epsilon\bm E^2-\frac{\bm B^2}{\mu}\right),
 \qquad
 c=\frac{1}{\sqrt{\mu\epsilon}}.
\end{equation}
The electric and magnetic fields are
\begin{equation}
 \bm E=-\frac{\partial\bm A}{\partial t}-\bm\nabla\phi,
 \qquad
 \bm B=\bm\nabla\times\bm A.
\end{equation}
They are invariant under the following gauge transformations
\begin{equation}
 \bm A\longmapsto\bm A+\bm\nabla\Theta,
 \qquad
 \phi\longmapsto\phi-\frac{\partial\Theta}{\partial t}.
\end{equation}
Below we choose the Weyl and Coulomb gauges:
\begin{equation}
 \phi=0,
 \qquad
 \bm\nabla\cdot\bm A=0.
\end{equation}

\subsection{One-dimensional optical cavity}\label{section:1d optical cavity}

In a cavity which is confined along $x,z$ direction but has length $L$ along $y$ direction, we consider one linearly polarized mode propagating along $y$,
\begin{equation}
 \bm A=(A(y,t),0,0),
 \qquad
 \bm E=(E_x,0,0),
 \qquad
 \bm B=(0,0,B_z),
\end{equation}
with
\begin{equation}
 E_x=-\frac{\partial A}{\partial t},
 \qquad
 B_z=-\frac{\partial A}{\partial y}.
\end{equation}
The cavity Lagrangian is
\begin{equation}
 L=\int_0^L\dif y\,
 \left[
 \frac{\epsilon}{2}\left(\frac{\partial A}{\partial t}\right)^2
 -\frac{1}{2\mu}\left(\frac{\partial A}{\partial y}\right)^2
 \right].
\end{equation}


\subsubsection{Standing wave boundary conditions}
Impose perfectly reflecting standing-wave boundary conditions
\begin{equation}
 \boxed{A(0,t)=A(L,t)=0.}
\end{equation}
The normalized spatial modes are
\begin{equation}
 u_n(y)=\sqrt{\frac{2}{L}}\sin(k_ny),
 \qquad
 k_n=\frac{n\pi}{L},
 \qquad
 n=1,2,3,\ldots,
\end{equation}
and satisfy
\begin{equation}
 \int_0^L\dif y\,u_n(y)u_m(y)=\delta_{nm}.
\end{equation}
Expand the field as
\begin{equation}
 \boxed{A(y,t)=\sum_{n=1}^{\infty}q_n(t)u_n(y)
 =\sqrt{\frac{2}{L}}\sum_{n=1}^{\infty}
 q_n(t)\sin(k_ny).}
\end{equation}
Then
\begin{align}
 E_x(y,t)&=-\sqrt{\frac{2}{L}}
 \sum_{n=1}^{\infty}\dot q_n(t)\sin(k_ny),\\
 B_z(y,t)&=-\sqrt{\frac{2}{L}}
 \sum_{n=1}^{\infty}k_nq_n(t)\cos(k_ny).
\end{align}

\subsubsection{Mode Lagrangian}
Using
\begin{align}
 \int_0^L\dif y\,\sin(k_ny)\sin(k_my)&=\frac{L}{2}\delta_{nm},\\
 \int_0^L\dif y\,\cos(k_ny)\cos(k_my)&=\frac{L}{2}\delta_{nm},
\end{align}
one obtains the following Lagrangian
\begin{equation}
 \boxed{L=\sum_{n=1}^{\infty}
 \left(\frac{\epsilon}{2}\dot q_n^{\,2}
 -\frac{k_n^2}{2\mu}q_n^2\right).}
\end{equation}
Thus there is one real harmonic oscillator for each standing-wave mode.

\subsection{Canonical Quantization}

The momentum conjugate to $q_n$ is
\begin{equation}
 p_n=\frac{\partial L}{\partial\dot q_n}=\epsilon\dot q_n.
\end{equation}
Canonical quantization gives
\begin{equation}
 \boxed{\comm{\hat q_n}{\hat p_m}=\imth\hbar\delta_{nm},}
 \qquad
 \comm{\hat q_n}{\hat q_m}=0,
 \qquad
 \comm{\hat p_n}{\hat p_m}=0.
\end{equation}
The Hamiltonian is
\begin{align}
 \hat H
 &=\sum_{n=1}^{\infty}
 \left(\hat p_n\dot{\hat q}_n-L_n\right)\\
 &=\boxed{\sum_{n=1}^{\infty}
 \left(\frac{\hat p_n^2}{2\epsilon}
 +\frac{k_n^2}{2\mu}\hat q_n^2\right)}.
\end{align}
Identifying each term with
\begin{equation}
 \frac{\hat p_n^2}{2M}+\frac12M\omega_n^2\hat q_n^2
\end{equation}
gives
\begin{equation}
 M=\epsilon,
 \qquad
 \omega_n^2=\frac{k_n^2}{\mu\epsilon}=c^2k_n^2.
\end{equation}
Therefore,
\begin{equation}
 \boxed{\omega_n=ck_n=\frac{n\pi c}{L}.}
\end{equation}

\subsubsection{Creation and annihilation operators}
Define
\begin{align}
 \hat a_n
 &=\sqrt{\frac{\epsilon\omega_n}{2\hbar}}\hat q_n
 +\frac{\imth}{\sqrt{2\hbar\epsilon\omega_n}}\hat p_n,\\
 \hat a_n^\dagger
 &=\sqrt{\frac{\epsilon\omega_n}{2\hbar}}\hat q_n
 -\frac{\imth}{\sqrt{2\hbar\epsilon\omega_n}}\hat p_n.
\end{align}
Then
\begin{equation}
 \comm{\hat a_n}{\hat a_m^\dagger}=\delta_{nm}.
\end{equation}
The inverse relations are
\begin{align}
 \hat q_n&=\sqrt{\frac{\hbar}{2\epsilon\omega_n}}
 (\hat a_n+\hat a_n^\dagger),\\
 \hat p_n&=-\imth\sqrt{\frac{\hbar\epsilon\omega_n}{2}}
 (\hat a_n-\hat a_n^\dagger).
\end{align}
Thus
\begin{equation}
 \boxed{\hat H=\sum_{n=1}^{\infty}
 \hbar\omega_n\left(\hat a_n^\dagger\hat a_n+\frac12\right).}
\end{equation}
The total zero-point energy is
\begin{equation}
 E_{0}=\frac12\sum_{n=1}^{\infty}\hbar\omega_n.
\end{equation}

\subsubsection{Photon number states}
For one mode,
\begin{equation}
 \hat N_n=\hat a_n^\dagger\hat a_n,
 \qquad
 \hat N_n\dket{N_n}=N_n\dket{N_n}.
\end{equation}
Each excitation carries energy $\hbar\omega_n$, and
\begin{equation}
 E_n(N_n)=\hbar\omega_n\left(N_n+\frac12\right).
\end{equation}

\subsection{Quantized Standing-Wave Fields}

\subsubsection{Vector potential}
In the Schr\"odinger picture,
\begin{equation}
 \boxed{
 \hat A(y)=\sum_{n=1}^{\infty}
 \sqrt{\frac{\hbar}{\epsilon L\omega_n}}
 \sin(k_ny)(\hat a_n+\hat a_n^\dagger).}
\end{equation}
In the Heisenberg picture,
\begin{equation}
 \boxed{
 \hat A(y,t)=\sum_{n=1}^{\infty}
 \sqrt{\frac{\hbar}{\epsilon L\omega_n}}
 \sin(k_ny)
 \left(\hat a_ne^{-\imth\omega_nt}
 +\hat a_n^\dagger e^{\imth\omega_nt}\right).}
\end{equation}
The boundary conditions are automatically satisfied because
\begin{equation}
 \sin(k_n0)=\sin(k_nL)=0.
\end{equation}

\subsubsection{Electric and magnetic fields}
Using $\hat E_x=-\partial_t\hat A$,
\begin{equation}\label{electric field}
 \boxed{
 \hat E_x(y,t)=\imth\sum_{n=1}^{\infty}
 \sqrt{\frac{\hbar\omega_n}{\epsilon L}}
 \sin(k_ny)
 \left(\hat a_ne^{-\imth\omega_nt}
 -\hat a_n^\dagger e^{\imth\omega_nt}\right).}
\end{equation}
Using $\hat B_z=-\partial_y\hat A$,
\begin{equation}
 \boxed{
 \hat B_z(y,t)=-\sum_{n=1}^{\infty}
 k_n\sqrt{\frac{\hbar}{\epsilon L\omega_n}}
 \cos(k_ny)
 \left(\hat a_ne^{-\imth\omega_nt}
 +\hat a_n^\dagger e^{\imth\omega_nt}\right).}
\end{equation}

